% GENERATED FILE. Edit canonical lessons, then run npm run generate:books.
% canonical-source-hash: fnv1a64:a36431a12549df0d
% canonical-lessons: FA-C290-miravam, FA-C290-raftam, FA-C290-gozashte, FA-W290-zal, FA-C290-miraftam, FA-R290-the-past-two-ways

\chapter{The Past, Seen Two Ways}
\label{ch:fa-the-past-seen-two-ways}
\hlchaptermodality{290}

\begin{chapteropening}
\textbf{By the end of this chapter:} \emph{I can say what I do now, what I did, and what was going on, setting a scene in the mi- past and dropping events into it with the plain past.}

Complete FA-R290-the-past-two-ways: a scene in the mi- past and an event that cuts into it, joined with dar hali ke.
\end{chapteropening}

\section[\texorpdfstring{\fa{میروم} (miravam)}{میروم (miravam)}]{\fa{میروم} (miravam) — the present tense, built from the stem you stored}

\label{lesson:FA-C290-miravam}



\begin{quote}
Give the pair for \textquotedblleft{}to go\textquotedblright{}: the infinitive, then its present stem. (\emph{Raftan}, \emph{rav-}.) Then the present stem of \emph{kardan}. (\emph{Kon-}.)

Chapter six asked you to store those stems before there was a tense to run them on. This lesson is the tense.
\end{quote}

\begin{grammarlens}[title={Grammar Lens: three pieces, in one order}]
A Persian present is three pieces: \textbf{mi-}, then the present stem, then an ending for the person.

\begin{quote}
mi- + rav + -am $\to$ \textbf{\fa{میروم}} — \emph{miravam} — I go, I'm going
\end{quote}

The endings are the six chapter seven named: \emph{-am, -i, -ad, -im, -id, -and}. So: \emph{miravam}, I go; \emph{miravi}, you go; \emph{miravad}, he or she goes; \emph{miravim}, we go; \emph{miravid}, you go, to several people or politely; \emph{miravand}, they go.

The pronouns are \emph{man}, \emph{to}, \emph{u} for \textquotedblleft{}he\textquotedblright{} or \textquotedblleft{}she\textquotedblright{}, \emph{mâ} for \textquotedblleft{}we\textquotedblright{}, \emph{shomâ} and \emph{ânhâ}. The ending already says who, so \textbf{\fa{میروم}} is a whole sentence.

Any stored stem works the same way: \emph{mikonam}, I do; \emph{mikhânam}, I read; \emph{midânam}, I know. A compound verb keeps its noun in front, \textbf{\fa{فکر} \fa{میکنم}}, \emph{fekr mikonam}, \textquotedblleft{}I think\textquotedblright{}.

One tense covers \textquotedblleft{}I go\textquotedblright{}, \textquotedblleft{}I'm going\textquotedblright{} and often \textquotedblleft{}I'll go\textquotedblright{}: \textbf{\fa{فردا} \fa{میروم}} — \emph{fardâ miravam} — I'm going tomorrow.

Printed Persian often sets \textbf{\fa{می}} off with a narrow break. This book writes it joined to the verb, \textbf{\fa{میروم}}, and it is read the same either way.
\end{grammarlens}

\begin{grammarlens}[title={Grammar Lens: ne-, and the verb that refuses mi-}]
\textquotedblleft{}Not\textquotedblright{} is \textbf{ne-}, in front of everything: \textbf{\fa{نمیروم}}, \emph{nemiravam}, I'm not going; \textbf{\fa{نمیدانم}}, \emph{nemidânam}, I don't know.

\emph{Dâshtan}, \textquotedblleft{}to have\textquotedblright{}, takes no \emph{mi-} in this tense: \textbf{\fa{دارم}}, \emph{dâram}, I have, and \emph{dust dâram}, I like. Its \textquotedblleft{}not\textquotedblright{} is \emph{na-}: \emph{nadâram}, \emph{dust nadâram}. A few compounds on \emph{dâshtan} do take it, such as \emph{negah dâshtan}, \textquotedblleft{}to keep\textquotedblright{}: \emph{negah midâram}.
\end{grammarlens}

\begin{cousinweb}
Histories of Persian usually trace \emph{mi-} to an older adverb, \emph{hamē}, \textquotedblleft{}continually\textquotedblright{}, once a word of its own in front of the verb. EtymDB has no entry for the prefix, so this is the usual account rather than a sourced one.
\end{cousinweb}

\subsection*{Your turn}
Say these aloud:

\begin{itemize}
\raggedright
  \item \emph{miravam, miravi, miravad}, then \emph{miravim, miravid, miravand}
  \item \textquotedblleft{}I know\textquotedblright{}, then \textquotedblleft{}I don't know\textquotedblright{} — \emph{midânam}, \emph{nemidânam}
  \item \textquotedblleft{}I'm going tomorrow\textquotedblright{} — \emph{fardâ miravam}
  \item \textquotedblleft{}I like\textquotedblright{}, then \textquotedblleft{}I don't like\textquotedblright{} — \emph{dust dâram}, \emph{dust nadâram}
\end{itemize}

\begin{tcolorbox}[breakable,colback=teal!4,colframe=teal!35!black,title={Before you move on}]
What are the three pieces of a present? (\emph{Mi-}, the present stem, an ending.) \textquotedblleft{}We go\textquotedblright{}? (\emph{\textbf{Miravim.}}) \textquotedblleft{}I think\textquotedblright{}? (\emph{\textbf{Fekr mikonam.}}) \textquotedblleft{}I'm not going\textquotedblright{}? (\emph{\textbf{Nemiravam.}}) Which verb leaves \emph{mi-} off? (\emph{Dâshtan}: \emph{\textbf{dâram}}.)
\end{tcolorbox}
\section[\texorpdfstring{\fa{رفتم} (raftam)}{رفتم (raftam)}]{\fa{رفتم} (raftam) — the past, and the mark on an object you mean}

\label{lesson:FA-C290-raftam}



\begin{quote}
Say \textquotedblleft{}I'm going\textquotedblright{}, then \textquotedblleft{}I'm not going\textquotedblright{}. (\emph{Miravam}, \emph{nemiravam}.)

Now move it to yesterday. The past needs no new stem at all: the infinitive already holds it.
\end{quote}

\begin{grammarlens}[title={Grammar Lens: the past stem, and the same endings}]
Take an infinitive and drop its \emph{-an}: \emph{raftan} gives \textbf{\fa{رفت}}, \emph{raft}, the past stem chapter six named. Add the endings you know, with one change: \textquotedblleft{}he\textquotedblright{} or \textquotedblleft{}she\textquotedblright{} takes no ending at all.

\begin{quote}
\textbf{\fa{رفتم}} — \emph{raftam}, I went · \emph{rafti}, you went · \emph{raft}, she went · \emph{raftim}, we went · \emph{raftid}, you went · \emph{raftand}, they went
\end{quote}

Every verb gives up its past stem this way, \emph{-tan} and \emph{-dan} alike: \emph{budan}, \emph{budam}, I was, and \emph{bud}, it was; \emph{khordan}, \emph{khordam}, I ate; \emph{didan}, \emph{didam}, I saw. \textquotedblleft{}Not\textquotedblright{} here is \emph{na-}: \emph{naraftam}, I didn't go.

\begin{quote}
\textbf{\fa{دیروز} \fa{به} \fa{خانه} \fa{رفتم}.} — \emph{diruz be khâne raftam} — Yesterday I went home.
\end{quote}

\emph{Be} is \textquotedblleft{}to\textquotedblright{}. The verb comes last, and the time usually comes first.
\end{grammarlens}

\begin{grammarlens}[title={Grammar Lens: râ, for the one you mean}]
\begin{quote}
\textbf{\fa{فیلم} \fa{را} \fa{دیدم}.} — \emph{film râ didam} — I saw the film.
\end{quote}

\textbf{\fa{را}}, \emph{râ}, follows an object that is a particular one, the film both of you have in mind. Leave it out and the object is any film: \emph{film didam}, \textquotedblleft{}I watched a film\textquotedblright{}. \textbf{\fa{کتاب} \fa{را} \fa{خواندم}}, \emph{ketâb râ khândam}: I read the book.
\end{grammarlens}

\begin{cousinweb}
Grammars of Persian usually trace the \emph{-t} or \emph{-d} at the end of a past stem to an old participle ending, so \emph{raft} first meant \textquotedblleft{}gone\textquotedblright{} and \emph{kard}, \textquotedblleft{}done\textquotedblright{}. Old Persian \emph{kṛta}, \textquotedblleft{}done\textquotedblright{}, is usually given as the older shape of \emph{kard}. That is why the stem stands alone for \textquotedblleft{}she went\textquotedblright{}: \emph{raft} is \textquotedblleft{}gone\textquotedblright{} with nobody added.
\end{cousinweb}

\subsection*{Your turn}
Say these aloud:

\begin{itemize}
\raggedright
  \item \emph{raftam, rafti, raft}, then \emph{raftim, raftid, raftand}
  \item \textquotedblleft{}Yesterday I went home\textquotedblright{} — \emph{diruz be khâne raftam}
  \item \textquotedblleft{}I saw the film\textquotedblright{}, then its \textquotedblleft{}not\textquotedblright{} — \emph{film râ didam}, \emph{film râ nadidam}
  \item \textquotedblleft{}I was\textquotedblright{}, then \textquotedblleft{}it was\textquotedblright{} — \emph{budam}, \emph{bud}
\end{itemize}

\begin{tcolorbox}[breakable,colback=teal!4,colframe=teal!35!black,title={Before you move on}]
How do you find a past stem? (Drop the \emph{-an} of the infinitive.) \textquotedblleft{}I ate\textquotedblright{}? (\emph{\textbf{Khordam.}}) Which person takes no ending? (\textquotedblleft{}He\textquotedblright{} or \textquotedblleft{}she\textquotedblright{}: \emph{\textbf{raft}}.) What does \emph{râ} tell you? (That the object is \textbf{a particular one}.)
\end{tcolorbox}
\section[\texorpdfstring{\fa{گذشته} (gozashte)}{گذشته (gozashte)}]{\fa{گذشته} (gozashte) — last year, last month, and the past itself}

\label{lesson:FA-C290-gozashte}



\begin{quote}
Say \textquotedblleft{}I went\textquotedblright{}, then \textquotedblleft{}we went\textquotedblright{}. (\emph{Raftam}, \emph{raftim}.) Then the one word for \textquotedblleft{}last year\textquotedblright{}. (\emph{Pârsâl}.)

\emph{Pârsâl} works only for the year. For last month and last week, Persian has one word that fits them all.
\end{quote}

\subsection*{You'll want to know: \fa{گذشته}}
\textbf{\fa{گذشته}} — \emph{gozashte} — \textquotedblleft{}last\textquotedblright{}, of a stretch of time, and \textquotedblleft{}the past\textquotedblright{}.

It follows its noun with the ezafe, the same \emph{-e} as in \emph{esm-e man}:

\begin{quote}
\textbf{\fa{سال} \fa{گذشته} \fa{به} \fa{تهران} \fa{رفتم}.} — \emph{sâl-e gozashte be Tehrân raftam} — Last year I went to Tehran.
\end{quote}

So \emph{mâh-e gozashte} is \textquotedblleft{}last month\textquotedblright{}, and \emph{hafte-ye gozashte} is \textquotedblleft{}last week\textquotedblright{}: after a final vowel the ezafe is said \emph{-ye}, and in plain print it is often not written at all, \textbf{\fa{هفته} \fa{گذشته}}.

Standing alone it is a noun: \textbf{\fa{در} \fa{گذشته}}, \emph{dar gozashte}, \textquotedblleft{}in the past\textquotedblright{}.

Say it in three beats, \emph{go-zash-te}. Its second letter is one you have not written yet. The next lesson takes it out and puts it in your hand.

\begin{cousinweb}
\emph{gozashte} is a past stem with \emph{-e} on the end, the way Persian makes a participle: \textbf{\fa{گذشت}}, \emph{gozasht}, from \emph{gozashtan}, \textquotedblleft{}to pass\textquotedblright{}, and then \emph{-e}: \textquotedblleft{}passed\textquotedblright{}, the time that has gone by. \emph{Rafte}, \textquotedblleft{}gone\textquotedblright{}, is built the same way from \emph{raft}.

EtymDB takes \emph{gozashtan} back to a Middle Persian verb, written \emph{wtštn'} in Pahlavi letters, \textquotedblleft{}to pass, to cross\textquotedblright{}. It records no older parent, so no English cousin is claimed here.
\end{cousinweb}

\subsection*{Your turn}
Say these aloud:

\begin{itemize}
\raggedright
  \item \emph{gozashte}, slowly, then at speed — \emph{go-zash-te}
  \item \textquotedblleft{}Last year I went to Tehran\textquotedblright{} — \emph{sâl-e gozashte be Tehrân raftam}
  \item \textquotedblleft{}last month\textquotedblright{}, then \textquotedblleft{}last week\textquotedblright{} — \emph{mâh-e gozashte}, \emph{hafte-ye gozashte}
  \item \textquotedblleft{}in the past\textquotedblright{} — \emph{dar gozashte}
\end{itemize}

\begin{tcolorbox}[breakable,colback=teal!4,colframe=teal!35!black,title={Before you move on}]
\textquotedblleft{}Last month\textquotedblright{}? (\emph{\textbf{Mâh-e gozashte.}}) What does \emph{gozashte} mean, taken apart? (\textbf{Passed}: the past stem \emph{gozasht} with \emph{-e}.) Which word still does the same job for the year alone? (\emph{\textbf{Pârsâl.}})
\end{tcolorbox}
\section[\texorpdfstring{\fa{ذ} (zâl)}{ذ (zâl)}]{\fa{ذ} (zâl) — the letter inside \fa{گذشته}}

\label{lesson:FA-W290-zal}



\begin{quote}
Before the new one: write \fa{د} and \fa{ز} once each, and name them.

Now say \textbf{\fa{گذشته}} (\emph{gozashte}), \textquotedblleft{}last\textquotedblright{}, as in \emph{sâl-e gozashte}. You met it a moment ago. This lesson takes one letter out of it and puts it in your hand.
\end{quote}

\subsection*{Script you'll notice: \fa{ذ}}
\textbf{\fa{ذ}} — \emph{zâl}, the sound \emph{z}.

It is the second letter of \textbf{\fa{گذشته}}, straight after \textbf{\fa{گ}}. It is \textbf{\fa{د}} with \textbf{one dot above}, and like \textbf{\fa{د}} it takes the hand of the letter on its right and offers nothing to the left: after \textbf{\fa{ذ}} the pen lifts, and \textbf{\fa{شته}} starts fresh.

You now hold three ways to write \emph{z}: \textbf{\fa{ز}}, \textbf{\fa{ظ}}, and \textbf{\fa{ذ}}. Persian says them the same. Most words with \textbf{\fa{ذ}} came from Arabic; \textbf{\fa{گذشته}} is one of a few Persian words that kept it. The usual explanation is that early New Persian softened a \emph{d} after a vowel towards the \emph{th} of English \emph{this}, and scribes wrote that sound with \textbf{\fa{ذ}}.

\subsection*{Writing: \fa{ذ}}
Put your pen on \fa{ذ} and follow its line. Copy the shape you can see — slowly, and larger than it is printed.

\begin{quote}
This book does not yet tell you \textbf{where to start this letter or which way to travel}. It is not written down here until it can be written down with a source. Copying what is in front of you needs no such source, and it is how the shape gets into your hand in the meantime.
\end{quote}

\subsection*{Your turn}
\begin{itemize}
\raggedright
  \item \emph{Find:} \fa{ذ} inside \textbf{\fa{گذشته}}, and point at it
  \item \emph{Trace it:} \fa{ذ} three times, saying \emph{z} as you finish each one
  \item \emph{Write it:} \fa{د}, then \fa{ذ}, from memory
  \item \emph{Write it:} \textbf{\fa{گذشته}}, with the break after \fa{ذ}
  \item \emph{Read it:} \textbf{\fa{گذشته}} aloud once more
\end{itemize}

\begin{tcolorbox}[breakable,colback=teal!4,colframe=teal!35!black,title={Before you move on}]
Which letter is this — \fa{ذ}? (\textbf{zâl}, \emph{z}.) Which word did it come from? (\textbf{\fa{گذشته}}, \emph{gozashte}, \textquotedblleft{}last\textquotedblright{}.) What is the difference between \fa{د} and \fa{ذ}? (\textbf{One dot above.})
\end{tcolorbox}
\section[\texorpdfstring{\fa{میرفتم} (miraftam)}{میرفتم (miraftam)}]{\fa{میرفتم} (miraftam) — the past that is still running}

\label{lesson:FA-C290-miraftam}



\begin{quote}
Say \textquotedblleft{}I'm going\textquotedblright{}, then \textquotedblleft{}I went\textquotedblright{}. (\emph{Miravam}, \emph{raftam}.)

\emph{Raftam} has an edge: the going happened, and it is over. Persian has a second past with no edge on it, and a story needs both.
\end{quote}

\begin{grammarlens}[title={Grammar Lens: mi- on the past stem}]
Put the \emph{mi-} of the present onto the past stem, and keep the past endings:

\begin{quote}
mi- + raft + -am $\to$ \textbf{\fa{میرفتم}} — \emph{miraftam} — I was going; I used to go
\end{quote}

So \emph{mikhândam}, I was reading; \emph{mikhordam}, I was eating; \emph{nemiraftam}, I wasn't going. \textquotedblleft{}She\textquotedblright{} again takes no ending: \emph{miraft}, she was going.

It answers two English shapes. \textquotedblleft{}Was ‑ing\textquotedblright{}: the action was under way. \textquotedblleft{}Used to\textquotedblright{}: it happened again and again.

\begin{quote}
\textbf{\fa{هر} \fa{روز} \fa{به} \fa{مدرسه} \fa{میرفتم}.} — \emph{har ruz be madrese miraftam} — Every day I used to go to school.
\end{quote}

\emph{Budan} normally takes no \emph{mi-} in the past. \emph{Bud} is already the background: \textbf{\fa{هوا} \fa{سرد} \fa{بود}}, \emph{havâ sard bud}, it was cold.
\end{grammarlens}

\begin{grammarlens}[title={Grammar Lens: the scene, and the events}]
A told story splits its verbs in two. The \emph{mi-} past paints the \textbf{scene}: what was going on, what things were like. The plain past drops the \textbf{events} into it: what happened next.

\begin{quote}
\textbf{\fa{در} \fa{حالی} \fa{که} \fa{کتاب} \fa{میخواندم،} \fa{تلفن} \fa{زنگ} \fa{زد}.} — \emph{dar hâli ke ketâb mikhândam, telefon zang zad} — While I was reading a book, the phone rang.
\end{quote}

The reading is scene, still running when the ring cuts into it. The ring is an event, over in a moment. \emph{Dar hâli ke}, \textquotedblleft{}while\textquotedblright{}, is the joint between them.

\begin{quote}
\textbf{\fa{هوا} \fa{سرد} \fa{بود} \fa{و} \fa{خورشید} \fa{میتابید}.} — \emph{havâ sard bud va khorshid mitâbid} — It was cold, and the sun was shining.
\end{quote}

Two pieces of scene, and nothing has happened yet. That is how a story opens.
\end{grammarlens}

\subsection*{Your turn}
Say these aloud:

\begin{itemize}
\raggedright
  \item \emph{miraftam, mirafti, miraft} — then \emph{raftam}, the one with an edge
  \item \textquotedblleft{}Every day I used to go to school\textquotedblright{} — \emph{har ruz be madrese miraftam}
  \item \emph{dar hâli ke ketâb mikhândam, telefon zang zad}
  \item a scene for last year — \emph{sâl-e gozashte havâ sard bud}
\end{itemize}

\begin{tcolorbox}[breakable,colback=teal!4,colframe=teal!35!black,title={Before you move on}]
How do you build \textquotedblleft{}I was reading\textquotedblright{}? (\emph{Mi-} on the past stem: \emph{\textbf{mikhândam}}.) In the phone sentence, which verb is the scene? (\emph{\textbf{Mikhândam}}, still running.) Which is the event? (\emph{\textbf{Zang zad}}, over in a moment.) Why is there no \emph{mi-} in \emph{havâ sard bud}? (\emph{Budan} takes none: \emph{\textbf{bud}} is already background.)
\end{tcolorbox}
\section[(review)]{Review — the scene, and what cut into it}

\label{lesson:FA-R290-the-past-two-ways}



\begin{quote}
Say \textquotedblleft{}I read\textquotedblright{} for today, then for yesterday. (\emph{Mikhânam}, \emph{khândam}.) Then \textquotedblleft{}I was reading\textquotedblright{}. (\emph{Mikhândam}.)
\end{quote}

\begin{grammarlens}[title={What you've built}]
Four lessons, and three tenses from two stems. The present stem with \emph{mi-}: \emph{miravam}, I go. The past stem alone: \emph{raftam}, I went, an event with an edge. The past stem with \emph{mi-}: \emph{miraftam}, I was going, the scene. And \emph{ne-} or \emph{na-} in front for \textquotedblleft{}not\textquotedblright{}.

\emph{Gozashte} sets the time, and \emph{râ} marks the object you mean:

\begin{quote}
\textbf{\fa{هفته} \fa{گذشته} \fa{فیلم} \fa{را} \fa{دیدم}.} — \emph{hafte-ye gozashte film râ didam} — Last week I saw the film.
\end{quote}

And \emph{dar hâli ke} joins a scene to the event that cut into it:

\begin{quote}
\textbf{\fa{در} \fa{حالی} \fa{که} \fa{فیلم} \fa{را} \fa{میدیدم،} \fa{تلفن} \fa{زنگ} \fa{زد}.} — \emph{dar hâli ke film râ mididam, telefon zang zad} — While I was watching the film, the phone rang.
\end{quote}
\end{grammarlens}

\subsection*{Your turn}
Say these aloud:

\begin{itemize}
\raggedright
  \item a scene in two — \emph{havâ sard bud. ketâb mikhândam.}
  \item \textquotedblleft{}Last week I saw the film\textquotedblright{} — \emph{hafte-ye gozashte film râ didam}
  \item \emph{dar hâli ke film râ mididam, telefon zang zad}
  \item one of your own: \emph{dar hâli ke} with a \emph{mi-} past, then one thing that happened
\end{itemize}

\begin{tcolorbox}[breakable,colback=teal!4,colframe=teal!35!black,title={Before you move on}]
After \emph{dar hâli ke}, which past? (The \emph{\textbf{mi-}} past, the scene.) Which past has an edge? (The plain one: \emph{\textbf{didam}}.) \textquotedblleft{}Last month\textquotedblright{}? (\emph{\textbf{Mâh-e gozashte.}}) Which letter in \textbf{\fa{گذشته}} is new to your hand? (\textbf{\fa{ذ}}, \emph{zâl}.)
\end{tcolorbox}
